\section{Implications and Limits}
The same functional deletion mask can receive different task-level fault labels under two healthy-capable feedback laws. A neural threshold and a task requirement should therefore be reported as separate specifications, with the actuator schedule and task horizon attached. The two-sided activity control shows that the discrepancy is not explained solely by choosing an upper envelope for a deletion fault.

Evidence remains conditional on one v783 graph, one assumed visual map and a simple finite-horizon pursuit task. The laws share actuator limits but use different turning histories; relay performance does not establish lower control cost or a hardware protection policy. Geometry-informed forecasts have more information than neural counts, and false alarms remain substantial. Sparse failure clusters and fixed calibration limit uncertainty estimates. Functional row deletion is uncalibrated to physical bit errors, and no energy, area, real-time execution or biological fault tolerance is established.

\section{Conclusion}
An auditable paired assessment separates feedback-dependent task loss from neural-rate and activity criteria in a large connectome SNN. Independent evidence supports a 50\% deletion contrast at the specified one-second endpoint; the 25\% contrast remains uncertain and stopped-duration matching has little effect in the new cohort. The reusable result is a source-bound assessment and replay workflow with explicit task, timing and resource endpoints.
